\section{The bijection and its statistics}\label{sec:bijection}

Start with $\lambda^{(0)}=\lambda\in\Rk$ and define
\begin{equation}\label{eq:iterated-transfer}
 \lambda^{(i+1)}=\tau\bigl(\lambda^{(i)}\bigr)
 \quad\text{while }\lambda^{(i)}\ne\varnothing.
\end{equation}
The process terminates by \eqref{eq:size-transfer}. Let
\begin{equation}\label{eq:transfer-sequence}
 \mathbf{s}(\lambda)=(s_1,s_2,\ldots,s_m),
 \qquad s_i=\bk\bigl(\lambda^{(i-1)}\bigr),
\end{equation}
where $\lambda^{(m)}=\varnothing$. Define
\begin{equation}\label{eq:definition-phi}
 \Phi_k(\lambda)=\mathbf{s}(\lambda)'.
\end{equation}

\begin{proposition}\label{prop:sequence-properties}
Let $\lambda\in\Rk$ be nonempty. The sequence $\mathbf{s}(\lambda)$ is a
partition and satisfies
\begin{equation}\label{eq:gap-bound}
 0\leq s_i-s_{i+1}<k\quad(1\leq i<m),
 \qquad 0<s_m<k.
\end{equation}
Moreover, $\Phi_k(\lambda)$ is $k$-strict and satisfies the three identities
in \eqref{eq:statistic-preservation}.
\end{proposition}

\begin{proof}
Apply \Cref{lem:converse-transfer} at every step of
\eqref{eq:iterated-transfer}. It gives the first inequality in
\eqref{eq:gap-bound}. In the final nonempty partition all remaining parts lie
in a single block with top less than $k$, so $0<s_m<k$.

For a partition $\mathbf{s}$, the multiplicity of the part $i$ in its
conjugate is $s_i-s_{i+1}$, where $s_{m+1}=0$. Thus
\eqref{eq:gap-bound} implies that $\mathbf{s}'$ is $k$-strict. Telescoping
\eqref{eq:size-transfer} over the transfer sequence gives
\begin{equation}\label{eq:size-telescope}
 |\lambda|=s_1+s_2+\cdots+s_m=|\mathbf{s}'|.
\end{equation}
The length of $\mathbf{s}'$ is its largest part $s_1=\bk(\lambda)$.

It remains to identify the length statistic. A transfer changes no number of
parts except when a terminal block is present, in which case it deletes one
part. Such a block is present exactly when the corresponding $s_i$ is not a
multiple of $k$. Hence
\begin{equation}\label{eq:length-nonmultiples}
 \ell(\lambda)=\#\{i: k\nmid s_i\}.
\end{equation}
Let $\pi=\mathbf{s}'$. For every $j\geq0$,
\begin{equation}\label{eq:strip-count}
 \pi_{jk+1}-\pi_{(j+1)k}
 =\#\{i: jk<s_i<(j+1)k\}.
\end{equation}
Summing \eqref{eq:strip-count} and using \eqref{eq:k-alternating-length}
shows that its left-hand side is $\lk(\pi)$ and its right-hand side is the
quantity in \eqref{eq:length-nonmultiples}. This proves all three identities.
\end{proof}

\begin{proof}[Proof of \Cref{thm:main}]
The empty partition maps to itself. For a nonempty $\lambda$, the preceding
proposition shows that $\Phi_k$ has the required statistics. It remains to
prove that the map is a bijection on nonempty partitions. Let $\pi\in\Sk$ be
nonempty and put
$\mathbf{s}=\pi'=(s_1,\ldots,s_m)$. Since $\pi$ is $k$-strict,
\begin{equation}\label{eq:inverse-gap-bound}
 0\leq s_i-s_{i+1}<k\quad(1\leq i<m),
 \qquad 0<s_m<k.
\end{equation}
Set $\mu^{(m)}=\varnothing$. For $i=m,m-1,\ldots,1$, define recursively
\begin{equation}\label{eq:inverse-recursion}
 \mu^{(i-1)}=\sigma_{s_i}\bigl(\mu^{(i)}\bigr).
\end{equation}
Inductively, $\bk(\mu^{(i)})=s_{i+1}$, with $s_{m+1}=0$. Thus
\eqref{eq:inverse-gap-bound} is exactly the admissibility condition in
\Cref{lem:block-transfer}. That lemma gives
$\tau(\mu^{(i-1)})=\mu^{(i)}$ and
$\bk(\mu^{(i-1)})=s_i$. Therefore
$\mathbf{s}(\mu^{(0)})=\mathbf{s}$, so $\Phi_k(\mu^{(0)})=\pi$. The same
reverse recursion is forced by \Cref{lem:converse-transfer}; hence the
preimage is unique. This proves bijectivity and the theorem.
\end{proof}
